\newpage
\section{概述}

% 主题：AI+量子编译
% 分工：
% \begin{itemize}
%     \item 量子线路优化：陈铭瑜
%     \item 量子比特映射和路由：郑钊雨、潘豪
% \end{itemize}

在数字经济与前沿科技深度融合的时代背景下，人工智能(AI)和量子计算作为两大颠覆性技术，正从各自独立发展的阶段，迈向双向赋能、协同突破的关键时期。

近年来，量子计算和AI表现出了展现出了重构新时代的计算范式的可能性。同样作为计算范式，两者的协同作用成长为了一个新的研究方向\cite{carleo2019machine}，有文章将这一新的研究领域称作QAI\cite{acampora2025quantum}。
% 许多基于AI的技术驱动了量子计算中关键研究的发展。从NISQ(Noisy Intermediate-scale quantum)时代早期的优化量子控制和量子错误减轻策略，到设计新的量子算法。
这个领域被分为AI赋能的量子计算以及量子计算强化的AI。一方面，量子计算被期望能够加速经典ML模型的训练过程，利用NISQ设备上的变分算法，以及未来的容错量子计算机。
另一方面，强化学习能够改善量子的计算瓶颈以及逐渐增长的训练时间\cite{meyer2022survey}。而AI赋能的量子计算覆盖的领域更广泛，包括量子纠错\cite{overwater2022neural}、量子线路合成等。

本工作研究的主题是AI驱动的量子编译与线路优化。量子编译是让实现量子算法的高级量子程序在底层量子设备上运行的必要条件。
AI赋能有望通过先进人工智能技术，重构量子线路编译流程。通过AI对海量线路结构学习，探索出相较于经典方法更好的量子线路编译方案，提高量子线路在量子设备上的执行效率，降低量子线路执行时的噪声。目前，已经有许多AI驱动的量子编译工作，例如量子线路优化\cite{li2024quarl}、量子比特映射\cite{acampora2021deep,li2023quantum}、量子比特路由等。

为此，我们提出了一个新的编译框架，包含AI驱动的量子线路优化模块和量子比特映射及路由模块，能够减少量子线路的规模，并提升量子线路在具体物理设备，并在真实设备上得到了验证了框架的效果。